%%
%% Ergänzung zu Kapitel 2: Rechtliche Rahmenbedingungen
%%

\section{Pseudonymisierte Datenbestände}

Die wissenschaftlichen Selbstberichte sind nicht anonym, sondern pseudonymisiert. Für das intraindividuelle Studiendesign müssen die zwanzig Berichte derselben App-Installation zusammengeführt werden können. Zu diesem Zweck berechnet der Server aus dem lokal gespeicherten App-Token eine stabile Teilnehmerkennung. Die Berichtstabelle enthält damit zwar keine unmittelbar identifizierenden Angaben wie Name, E-Mail-Adresse oder Bankverbindung, aber eine wiederkehrende Kennung, anhand derer die Berichte derselben Person beziehungsweise Installation verknüpft werden können. Die Selbstberichte bleiben deshalb personenbezogene und wegen ihres Inhalts gesundheitsbezogene Daten.

Die Implementierung trennt drei Datenzusammenhänge. Die Registrierungsdatenbank enthält die für Anmeldung, Teilnahmeprüfung und Abrechnung erforderlichen direkten Identifikatoren. Die Studiendatenbank enthält die pseudonymisierten Selbstberichte, die experimentelle Bedingung und den Craving-Wert. Die Android-App hält das App-Token und den lokalen Studienfortschritt. Eine unmittelbare Fremdschlüsselbeziehung zwischen Registrierung und Selbstberichten wird nicht angelegt. Die Verbindung entsteht nur mittelbar über geheimnisgebundene Hashwerte, die in unterschiedlichen fachlichen Kontexten berechnet werden.

Für diese Hashwerte wird HMAC-SHA-256 mit einem ausschließlich serverseitig vorgehaltenen Geheimnis eingesetzt. Zusätzlich wird eine Domänentrennung verwendet. Aus demselben App-Token entstehen dadurch unterschiedliche Werte für die Zuordnung zur Registrierung, die Freigabe wissenschaftlicher Übertragungen und die pseudonyme Teilnehmerkennung. Verschiedene HMAC-Kontexte verhindern, dass ein in einer Tabelle gespeicherter Hash ohne Kenntnis des Geheimnisses unmittelbar mit einem Hash aus einer anderen Tabelle verglichen werden kann. Diese Maßnahme reduziert die Verknüpfbarkeit der Datenbestände gegenüber einem einfachen SHA-256-Hash des Tokens.

\section{Aktivierung und Datenminimierung}

Die Aktivierung der App erfolgt in zwei technischen Schritten. Zunächst übermittelt die App die normalisierte E-Mail-Adresse. Der Server prüft, ob die Registrierung per Double-Opt-In bestätigt wurde, die Studieninformation akzeptiert wurde und noch kein App-Token ausgegeben worden ist. Bei erfolgreicher Prüfung erzeugt er ein zufälliges UUID-v4-Token und gibt es an die App zurück. In der Registrierung wird zu diesem Zeitpunkt nur ein zeitlich befristeter, an E-Mail-Adresse und Token gebundener HMAC gespeichert. Das Token selbst wird nicht dauerhaft in der Datenbank abgelegt.

Anschließend bestätigt die App den Empfang des Tokens in einem zweiten Request. Erst nach dieser Bestätigung werden der temporäre Aktivierungshash und seine Gültigkeitsfrist entfernt, der Ausgabezeitpunkt vermerkt und die für spätere Statusprüfungen benötigten domänengetrennten Hashwerte gespeichert. Die Zweiteilung soll vermeiden, dass eine Registrierung bereits als aktiviert gilt, obwohl das Token die App wegen eines Verbindungsabbruchs nie erreicht hat. Die Gültigkeit des ersten Schritts ist auf fünf Minuten beschränkt; ein neuer Aktivierungsversuch ersetzt einen noch ausstehenden älteren Versuch.

Die E-Mail-Adresse wird ausschließlich für die Aktivierung verwendet. Wissenschaftliche Selbstberichte enthalten sie nicht. Auch der für Selbstberichte verwendete Endpunkt erhält nur das App-Token, den Craving-Wert und gegebenenfalls die App-Version. Situation und Versuchsbedingung werden serverseitig aus der Zahl bereits gespeicherter Berichte abgeleitet. Dadurch werden überflüssige, von der App manipulierbare Felder vermieden.

\section{Erneute Datenschutzzustimmung}

Die App unterstützt eine erneute Zustimmung, wenn der in der Registrierung gespeicherte Datenschutzstatus zurückgesetzt wurde. Der Status wird nicht anhand der E-Mail-Adresse, sondern anhand des domänengetrennten Registrierungs-Token-Hashs abgefragt. Der Server kann damit die zugehörige Registrierung finden, ohne das App-Token selbst zu speichern. Ein \texttt{POST}-Request liest den Status; ein \texttt{PUT}-Request mit dem ausdrücklich gesetzten Wert \texttt{dataprot=true} dokumentiert die Zustimmung und setzt den Zeitpunkt \path{dataprot_accepted_at}. Die Speicherung des Zeitpunkts ist idempotent: Eine wiederholte Übermittlung verändert den zuerst gespeicherten Annahmezeitpunkt nicht.

Bei fehlendem Token, unbekannter Registrierung, fehlerhafter Serverantwort oder nicht erreichbarem Dienst verhält sich die App grundsätzlich geschlossen. Die produktive Studienfunktion wird dann nicht geöffnet. Damit wird vermieden, dass ein technischer Fehler als Zustimmung interpretiert wird. Die Demo und das anonyme Feedback können getrennt davon angeboten werden, sofern sie keine wissenschaftlichen Selbstberichte übertragen.

\section{Datensparsame Benachrichtigungen}

Die App bietet allgemeine Informationsbenachrichtigungen und Erinnerungen an eine wieder verfügbare Studiensituation. Beide Funktionen sind freiwillig. Unter Android 13 oder neuer wird die Systemberechtigung \path{POST_NOTIFICATIONS} erst nach einer kontextbezogenen Erklärung angefordert. Wird die Berechtigung nicht erteilt oder später entzogen, bleiben Aktivierung, Demo, Feedback und die manuelle Durchführung der Studie möglich.

Benachrichtigungen dürfen keine E-Mail-Adresse, kein App-Token, keinen Rauchstatus und keinen Craving-Wert enthalten. Für Studienerinnerungen wird zusätzlich eine reduzierte öffentliche Version bereitgestellt, die auf dem Sperrbildschirm lediglich auf neue allgemeine Informationen hinweist. Erst nach dem Öffnen der App werden Feature-Freigabe, Aktivierungsstatus, ausstehende Übertragungen, Studienabschluss und Sperrzeit erneut geprüft. Eine Benachrichtigung stellt damit keine Autorisierung zum Start einer Studiensituation dar.

Die lokale Planung erfolgt über Android WorkManager. Das Betriebssystem darf die Ausführung wegen Energiesparmechanismen oder fehlender Netzverbindung verzögern. Die Erinnerung ist deshalb als unverbindliche Komfortfunktion und nicht als medizinisch oder studienmethodisch exakter Alarm einzuordnen.

\section{Lokale Verarbeitung}

Das App-Token wird mit AES-256-GCM verschlüsselt gespeichert. Der Schlüssel wird im Android Keystore erzeugt und verlässt den geschützten Schlüsselbereich nicht. Ist der Schlüssel verloren gegangen oder ungültig, bricht die App den Zugriff ab und verwendet den Wert nicht weiter.

Andere lokale Studienwerte, darunter Fortschritt, Sperrzeit, ein ausstehender Craving-Wert und der Abrechnungscode, liegen weiterhin in privaten App-Einstellungen. Die automatische Android-Sicherung ist deaktiviert, damit diese Werte nicht über allgemeine Cloud- oder Gerätebackups verbreitet werden.

\section{App-Verteilung}
Die derzeitige Verteilung erfolgt als direkt herunterladbares APK mit deutscher und englischer Installationsanleitung. Dieses Verfahren erlaubt die Durchführung außerhalb eines öffentlichen App-Stores, verlangt aber eine besonders klare Information über Herkunft, Version und Installation der Datei. Die Anleitung weist darauf hin, dass Dialoge je nach Browser und Android-Version abweichen können.
